\section{Persecutor, called apostle}\label{sec:paul-call}

Paul never allowed the Church to forget that he had persecuted it. Galatians
says he persecuted the assembly of God violently and tried to destroy it.
First Corinthians calls him least of the apostles and unfit for the name
because he persecuted the Church; Philippians places zeal and persecution
beside his former confidence in inherited standing. The admissions come from
Paul, not from a later enemy. They make grace, rather than heroic preparation,
the interpretive center of his vocation.

\subsection{What Paul says happened}

Paul does not narrate a road scene in his letters. He says that God, who had
set him apart before birth and called him through grace, was pleased to reveal
the Son to him so that he might proclaim him among the Gentiles (Gal
1:15--16). He asks the Corinthians whether he has not seen Jesus the Lord and
places himself last in a received list of Resurrection appearances (1 Cor
9:1; 15:3--11). The event was revelation, appearance, call, and commission.
Its reality grounded his apostleship; its graciousness exposed his former
violence.

The language echoes prophetic calls, especially a servant set apart from the
womb. Paul did not infer his mission from private spiritual achievement. Nor
did he claim that independence from human authorization meant independence
from the Church. The gospel and call came from Christ; Paul was baptized,
visited Cephas, met James, received and transmitted traditions, sought
fellowship with the pillars, and later organized a collection as a sign of
communion.

\subsection{Three accounts in Acts}

Acts tells the encounter three times: as third-person narrative in chapter 9,
as Paul's Temple-stair defense in chapter 22, and as his hearing before
Agrippa in chapter 26. All place it near Damascus amid persecuting activity.
A heavenly light, a fall to the ground, the voice of Jesus, and a commission
form the common center. The accounts vary in wording and narrative allocation.
The companions hear a voice in one formulation and do not hear or understand
it in another; they see light or no person; Ananias mediates healing and
commission in chapters 9 and 22, while chapter 26 places the mission directly
in Jesus' address.

These differences are not repaired by inventing a transcript. Luke retells
one event for different audiences and theological purposes. Acts 9 emphasizes
Paul's dependence on Ananias and the frightened Damascus community. Acts 22
locates his testimony within a Jewish defense and Ananias's devout observance
of the law. Acts 26 condenses the commission for a royal hearing. None of the
accounts mentions a horse. Later visual use of a horse was encountered only as
an uncollated art-historical lead.

Acts also sets Paul near Stephen's death, approving it and guarding garments,
then describes house-to-house arrests and authority sought from the high
priest. Paul's letters independently secure persecution but do not name
Stephen, Ananias, warrants, companions, blindness, or the Straight Street.
Those are elements of Acts' canonical testimony, not autobiographical details
that can be assigned the same source grade.

\subsection{Baptism, Arabia, and first contacts}

In Acts, Ananias lays hands on Saul; sight returns; Saul is baptized, eats, and
immediately proclaims Jesus in Damascus. Galatians instead stresses that Paul
did not immediately consult human beings or go to the earlier apostles. He
went to Arabia and returned to Damascus; only after three years did he go to
Jerusalem to visit Cephas, seeing also James. The accounts can coexist if Acts
compresses a longer period, but Acts never mentions Arabia, and Galatians does
not mention Ananias. The responsible biography preserves both.

``Arabia'' probably refers to Nabataean territory, but neither purpose nor
route is known. It may have involved mission, retreat, conflict, or several
activities. Second Corinthians 11:32--33 says the ethnarch of King Aretas
guarded Damascus and Paul escaped in a basket through a window. Acts also
narrates the escape but assigns the plot to opponents among the local Jews.
The difference warns against producing a single police report from two
theologically shaped memories.

\subsection{Conversion, call, and continuing identity}

Paul experienced a reversal profound enough to describe previous gains as
loss in comparison with knowing Christ. The persecutor became the
persecuted; the one who tried to destroy assemblies became responsible for
them. Yet the God who called him was the God of Israel, the gospel fulfilled
promises made beforehand, and the nations were grafted into Israel's olive
tree. Paul continued to call himself an Israelite, and Romans 9--11 joined the
mission to the nations with grief and hope for his kin.

\begin{studybox}{Grace remembered truthfully}
Paul neither hides persecution nor makes violence the source of his authority.
His apostolic self-understanding holds judgment and mercy together: what he
did was real, what he received was unearned, and grace became labor for
communities rather than a private exemption from accountability.
\end{studybox}
